\newacronym{ai}{AI}{Artificial Intelligence (slovensky „umelá inteligencia“)}
\newacronym{ann}{ANN}{Approximate Nearest Neighbor (slovensky „vyhľadávanie aproximovaných najbližších susedov“)}
\newacronym{api}{API}{Application Programming Interface (slovensky „aplikačné programové rozhranie“)}
\newacronym{asgi}{ASGI}{Asynchronous Server Gateway Interface (slovensky „asynchrónne rozhranie serverovej brány“)}
\newacronym{cli}{CLI}{Command Line Interface (slovensky „rozhranie príkazového riadku“)}
\newacronym{cpu}{CPU}{Central Processing Unit (slovensky „centrálna procesorová jednotka“)}
\newacronym{gpu}{GPU}{Graphics Processing Unit (slovensky „grafická procesorová jednotka“)}
\newacronym{hmi}{HMI}{Human-Machine Interface (slovensky „rozhranie človek-stroj“)}
\newacronym{hnsw}{HNSW}{Hierarchical Navigable Small World (slovensky „hierarchicky navigovateľné malé svety“)}
\newacronym{html}{HTML}{HyperText Markup Language (slovensky „hypertextový značkovací jazyk“)}
\newacronym{http}{HTTP}{HyperText Transfer Protocol (slovensky „hypertextový prenosový protokol“)}
\newacronym{iiot}{IIoT}{Industrial Internet of Things (slovensky „priemyselný internet vecí“)}
\newacronym{iot}{IoT}{Internet of Things (slovensky „internet vecí“)}
\newacronym{it}{IT}{Information Technology (slovensky „informačné technológie“)}
\newacronym{ivf}{IVF}{Inverted File Index (slovensky „index invertovaných súborov“)}
\newacronym{json}{JSON}{JavaScript Object Notation (slovensky „JavaScriptová objektová notácia“)}
\newacronym{kb}{KB}{Knowledge Base (slovensky „znalostná báza“)}
\newacronym{llm}{LLM}{Large Language Model (slovensky „veľký jazykový model“)}
\newacronym{mqtt}{MQTT}{Message Queuing Telemetry Transport (slovensky „protokol na prenos telemetrických správ“)}
\newacronym{nfc}{NFC}{Near Field Communication (slovensky „komunikácia v blízkom poli“)}
\newacronym{opcua}{OPC UA}{Open Platform Communications Unified Architecture (slovensky „zjednotená architektúra otvorenej platformovej komunikácie“)}
\newacronym{ot}{OT}{Operational Technology (slovensky „prevádzkové technológie“)}
\newacronym{pdf}{PDF}{Portable Document Format (slovensky „prenosný formát dokumentov“)}
\newacronym{plc}{PLC}{Programmable Logic Controller (slovensky „programovateľný logický automat“)}
\newacronym{qos}{QoS}{Quality of Service (slovensky „kvalita služieb“)}
\newacronym{qr}{QR}{Quick Response (slovensky „rýchla odozva“)}
\newacronym{rag}{RAG}{Retrieval-Augmented Generation (slovensky „generovanie rozšírené o vyhľadávanie“)}
\newacronym{ram}{RAM}{Random Access Memory (slovensky „pamäť s náhodným prístupom“)}
\newacronym{rest}{REST}{Representational State Transfer (slovensky „reprezentačný prenos stavu“)}
\newacronym{rlhf}{RLHF}{Reinforcement Learning from Human Feedback (slovensky „správaním posilňované učenie z ľudskej spätnej väzby“)}
\newacronym{rtu}{RTU}{Remote Terminal Unit (slovensky „vzdialená terminálová jednotka“)}
\newacronym{saas}{SaaS}{Software as a Service (slovensky „softvér ako služba“)}
\newacronym{scada}{SCADA}{Supervisory Control and Data Acquisition (slovensky „dohľadové riadenie a zber dát“)}
\newacronym{tcp}{TCP}{Transmission Control Protocol (slovensky „protokol riadenia prenosu“)}
\newacronym{ui}{UI}{User Interface (slovensky „používateľské rozhranie“)}
\newacronym{vram}{VRAM}{Video Random Access Memory (slovensky „video pamäť s náhodným prístupom“)}